\lecture{20}{PL-свойство композиций}

\sourcevideo
    {https://www.youtube.com/playlist?list=PLOzRYVm0a65fhKDS8cYIC7YCuVGJ4FOJx}
    {Week 6: Lecture 20: Advanced Results on PL Inequality: Part 2}

Лекция состоит из двух частей. Сначала доказывается, что композиция
сильно выпуклой функции с линейным отображением может потерять
сильную выпуклость, но сохраняет неравенство Поляка--Лоясиевича.
Затем исследуется робастность фиксированно-временного градиентного
потока к возмущению, исчезающему около множества решений.

\section{PL-условие и композиция}

Пусть дифференцируемая функция
\[
    f\colon\mathbb R^n\to\mathbb R
\]
достигает минимального значения
\[
    f^\star=\min_x f(x),
    \qquad
    \mathcal X^\star=\operatorname*{argmin}_x f(x).
\]

\begin{definition}[Неравенство Поляка--Лоясиевича]
Функция \(f\) удовлетворяет PL-неравенству с константой \(\mu>0\),
если
\[
    \frac12\lVert\nabla f(x)\rVert^2
    \ge
    \mu\bigl(f(x)-f^\star\bigr)
\]
для всех \(x\in\mathbb R^n\).
\end{definition}

PL-условие не требует выпуклости. Оно исключает неоптимальные
стационарные точки, поскольку
\[
    \nabla f(x)=0
    \quad\Longrightarrow\quad
    f(x)=f^\star,
\]
но не гарантирует единственности минимизатора.

Пусть
\[
    g\colon\mathbb R^m\to\mathbb R
\]
является \(\sigma\)-сильно выпуклой и
\[
    f(x)=g(Ax-b),
    \qquad
    A\in\mathbb R^{m\times n},
    \quad
    b\in\mathbb R^m.
\]
Тогда
\[
    \nabla f(x)=A^{\mathsf T}\nabla g(Ax-b).
\]
Из сильной выпуклости \(g\) следует
\[
\begin{aligned}
    f(y)
    \ge{}&
    f(x)
    +\langle\nabla f(x),y-x\rangle
    \\
    &+
    \frac\sigma2
    \lVert A(y-x)\rVert^2.
\end{aligned}
\]
Последнее слагаемое содержит \(\lVert A(y-x)\rVert\), поэтому
сильная выпуклость внешней функции не всегда переносится на
композицию.

Если \(A\) имеет полный столбцовый ранг, то
\[
    \lVert Az\rVert
    \ge
    \sigma_{\min}(A)\lVert z\rVert,
\]
и \(f\) является сильно выпуклой с константой
\[
    \sigma\sigma_{\min}(A)^2.
\]
Если же \(\ker A\ne\{0\}\), то
\[
    f(x+z)=f(x),
    \qquad
    z\in\ker A,
\]
поэтому глобальная сильная выпуклость невозможна.

\section{Сингулярное число, оценка Хоффмана и PL-свойство}

Пусть \(A\ne0\), но матрица может быть вырожденной. Обозначим её
наименьшее положительное сингулярное число через
\[
    \theta_A=\sigma_{\min}^{+}(A).
\]
Для каждого \(z\in(\ker A)^\perp\)
\[
    \lVert Az\rVert\ge\theta_A\lVert z\rVert.
\]
Сильная выпуклость \(g\) обеспечивает единственность решения задачи
\[
    \min_{u\in\operatorname{range}(A)-b}g(u).
\]
Обозначим его через \(u^\star\). Тогда
\[
    \mathcal X^\star
    =
    \{x\in\mathbb R^n:Ax-b=u^\star\}.
\]

Для заданного \(x\) выберем ближайший минимизатор
\[
    x^\star
    \in
    \operatorname*{argmin}_{z\in\mathcal X^\star}
    \lVert x-z\rVert.
\]
Тогда \(x-x^\star\in(\ker A)^\perp\), и потому
\[
    \lVert A(x-x^\star)\rVert
    \ge
    \theta_A\operatorname{dist}(x,\mathcal X^\star).
\]
Это частный случай оценки Хоффмана для линейной системы
\(Ax=Ax^\star\).

Подставляя \(y=x^\star\) в неравенство для композиции, получаем
\[
\begin{aligned}
    f^\star
    \ge{}&
    f(x)
    +\langle\nabla f(x),x^\star-x\rangle
    \\
    &+
    \frac{\sigma\theta_A^2}{2}
    \lVert x^\star-x\rVert^2.
\end{aligned}
\]
Это ограниченная сильная выпуклость относительно ближайшей точки
множества решений, а не обычная сильная выпуклость для произвольной
пары точек.

Положим
\[
    \mu_A=\sigma\theta_A^2.
\]
Для любых \(a,d\) выполняется
\[
    \langle a,d\rangle
    +\frac{\mu_A}{2}\lVert d\rVert^2
    \ge
    -\frac1{2\mu_A}\lVert a\rVert^2.
\]
Применяя это неравенство при
\[
    a=\nabla f(x),
    \qquad
    d=x^\star-x,
\]
получаем
\[
    f^\star
    \ge
    f(x)-\frac1{2\mu_A}\lVert\nabla f(x)\rVert^2.
\]
Следовательно,
\[
    \frac12\lVert\nabla f(x)\rVert^2
    \ge
    \mu_A\bigl(f(x)-f^\star\bigr).
\]

\begin{theorem}[PL-свойство линейной композиции]
Пусть \(g\) является \(\sigma\)-сильно выпуклой, \(A\ne0\), а
функция
\[
    f(x)=g(Ax-b)
\]
достигает минимума. Тогда \(f\) удовлетворяет PL-неравенству с
константой
\[
    \mu_A
    =
    \sigma
    \bigl(\sigma_{\min}^{+}(A)\bigr)^2.
\]
\end{theorem}

Матрица \(A\) может быть вырожденной. В этом случае функция не
является сильно выпуклой, но остаётся PL-функцией. При \(A=0\)
композиция постоянна и PL-неравенство выполняется тривиально.

\section{Наименьшие квадраты и границы переноса}

Для задачи
\[
    f(x)=\frac12\lVert Ax-b\rVert^2
\]
внешняя функция
\[
    g(u)=\frac12\lVert u\rVert^2
\]
является \(1\)-сильно выпуклой. Поэтому
\[
    \frac12\lVert\nabla f(x)\rVert^2
    \ge
    \bigl(\sigma_{\min}^{+}(A)\bigr)^2
    \bigl(f(x)-f^\star\bigr),
\]
где
\[
    \nabla f(x)=A^{\mathsf T}(Ax-b).
\]
Если \(A\) имеет полный столбцовый ранг, решение единственно и
функция сильно выпукла. При ненулевом ядре минимизаторы имеют вид
\[
    x^\star+z,
    \qquad
    z\in\ker A,
\]
но PL-свойство сохраняется. Поэтому градиентные методы могут линейно
сходиться по значению функции даже для вырожденной задачи наименьших
квадратов.

Тот же аргумент применим к любой композиции \(g(Ax-b)\), если
внешняя функция действительно сильно выпукла на рассматриваемой
области. Для стандартной нерегуляризованной логистической потери это
условие глобально обычно не выполнено; требуются дополнительные
ограничения или регуляризация.

\section{Возмущённый поток}

Пусть \(f\) удовлетворяет PL-неравенству с константой \(\mu>0\).
Сначала предположим, что минимизатор \(x^\star\) единственен. Выберем
\[
    p>2,
    \qquad
    1<q<2,
\]
и обозначим
\[
    \rho_p=\frac{p-2}{p-1},
    \qquad
    \rho_q=\frac{q-2}{q-1}.
\]
Рассмотрим возмущённую систему
\[
\begin{aligned}
    \dot x
    ={}&
    -c_1
    \frac{\nabla f(x)}{
        \lVert\nabla f(x)\rVert^{\rho_p}
    }
    \\
    &-c_2
    \frac{\nabla f(x)}{
        \lVert\nabla f(x)\rVert^{\rho_q}
    }
    +\varepsilon(x),
\end{aligned}
\]
где \(c_1,c_2>0\), а в стационарной точке правая часть определяется
равной нулю. Предполагается
\[
    \lVert\varepsilon(x)\rVert
    \le
    \ell\lVert x-x^\star\rVert^2.
\]
Отсюда \(\varepsilon(x^\star)=0\), поэтому оптимум сохраняется как
равновесие. При ненулевом остаточном возмущении точная сходимость к
\(x^\star\) в общем случае невозможна; можно ожидать лишь попадание
в окрестность.

Из предыдущей лекции PL-условие влечёт квадратичный рост:
\[
    f(x)-f^\star
    \ge
    \frac\mu2\lVert x-x^\star\rVert^2.
\]
Совместно с PL-неравенством это даёт
\[
    \lVert x-x^\star\rVert
    \le
    \frac1\mu\lVert\nabla f(x)\rVert.
\]
Следовательно,
\[
    \lVert\varepsilon(x)\rVert
    \le
    \frac\ell{\mu^2}
    \lVert\nabla f(x)\rVert^2.
\]
Обозначим
\[
    \overline\ell=\frac\ell{\mu^2},
    \qquad
    V(x)=f(x)-f^\star,
    \qquad
    s=\lVert\nabla f(x)\rVert.
\]
Вдоль возмущённого потока
\[
    \dot V
    \le
    -c_1s^{2-\rho_p}
    -c_2s^{2-\rho_q}
    +\overline\ell s^3.
\]
Положим
\[
    r_1=2-\rho_p=\frac{p}{p-1},
    \qquad
    r_2=2-\rho_q=\frac{q}{q-1}.
\]
Тогда
\[
    1<r_1<2,
    \qquad
    r_2>2.
\]

\section{Поглощение возмущения и робастная теорема}

При \(0\le s\le1\) из \(r_1<3\) следует
\[
    s^{r_1}\ge s^3.
\]
При \(s\ge1\) требуется
\[
    s^{r_2}\ge s^3,
\]
то есть \(r_2\ge3\). Это условие эквивалентно
\[
    1<q\le\frac32.
\]
Таким образом, для всех \(s\ge0\)
\[
    s^3\le s^{r_1}+s^{r_2},
\]
и потому
\[
    \dot V
    \le
    -\bigl(c_1-\overline\ell\bigr)s^{r_1}
    -\bigl(c_2-\overline\ell\bigr)s^{r_2}.
\]
Если
\[
    c_1>\overline\ell,
    \qquad
    c_2>\overline\ell,
\]
оба коэффициента положительны.

Введём
\[
    \alpha_1=\frac{r_1}{2}=\frac{p}{2(p-1)},
    \qquad
    \alpha_2=\frac{r_2}{2}=\frac{q}{2(q-1)}.
\]
Тогда
\[
    0<\alpha_1<1<\alpha_2.
\]
Из PL-неравенства \(s^2\ge2\mu V\) следует
\[
    \dot V
    \le
    -a_1V^{\alpha_1}
    -a_2V^{\alpha_2},
\]
где
\[
    a_1
    =
    \left(c_1-\frac\ell{\mu^2}\right)
    (2\mu)^{\alpha_1},
\]
\[
    a_2
    =
    \left(c_2-\frac\ell{\mu^2}\right)
    (2\mu)^{\alpha_2}.
\]

\begin{theorem}[Робастность к исчезающему возмущению]
Пусть \(f\) удовлетворяет PL-неравенству с константой \(\mu>0\) и
имеет единственный минимизатор \(x^\star\). Предположим, что
\[
    \lVert\varepsilon(x)\rVert
    \le
    \ell\lVert x-x^\star\rVert^2,
\]
а параметры удовлетворяют
\[
    p>2,
    \qquad
    1<q\le\frac32,
    \qquad
    c_1,c_2>\frac\ell{\mu^2}.
\]
Тогда \(x^\star\) является фиксированно-временно устойчивым
равновесием, причём
\[
    T(x_0)
    \le
    \frac1{a_1(1-\alpha_1)}
    +
    \frac1{a_2(\alpha_2-1)}.
\]
Правая часть не зависит от начальной точки.
\end{theorem}

Первое слагаемое потока поглощает возмущение при малом градиенте, а
второе --- при большом. Увеличение \(c_1,c_2\) уменьшает
теоретическую границу времени, но после дискретизации чрезмерное
усиление может ухудшить численную устойчивость.

\section{Несколько решений, дискретизация и дальнейший переход}

Если \(\mathcal X^\star\) содержит несколько точек, квадратичный рост
записывается как
\[
    f(x)-f^\star
    \ge
    \frac\mu2
    \operatorname{dist}(x,\mathcal X^\star)^2.
\]
Естественное условие на возмущение принимает вид
\[
    \lVert\varepsilon(x)\rVert
    \le
    \ell
    \operatorname{dist}(x,\mathcal X^\star)^2.
\]
Тогда исследуется фиксированно-временная устойчивость множества
решений, а не сходимость к заранее выбранному его элементу.

Явная дискретизация имеет вид
\[
\begin{aligned}
    x^{k+1}
    ={}&
    x^k
    -\eta c_1
    \frac{\nabla f(x^k)}{
        \lVert\nabla f(x^k)\rVert^{\rho_p}
    }
    \\
    &-\eta c_2
    \frac{\nabla f(x^k)}{
        \lVert\nabla f(x^k)\rVert^{\rho_q}
    }
    +\eta\varepsilon(x^k).
\end{aligned}
\]
Непрерывная фиксированно-временная гарантия не переносится на эту
схему автоматически. При слишком большом шаге или коэффициентах
возможны перескакивание через решение, колебания и расходимость. В
лекции утверждается существование достаточно малого шага, при котором
последовательность попадает в заданную окрестность, но универсальная
явная формула для общего нелинейного случая не приводится.

Далее курс переходит от задач без ограничений к постановкам
\[
    \min_x f(x)
    \qquad
    \text{при условии}
    \qquad
    Ax=b.
\]
Основными инструментами станут функция Лагранжа, дополненная функция
Лагранжа, метод множителей, двойственный подъём и ADMM.
